
on ait : @@M0@@.

\textbf{Lemme 2.11.} \emph{Soient @@M1@@ un schéma noethérien et @@M2@@ un @@M3@@ schéma projectif, muni d'un faisceau très ample @@M4@@. On suppose que les fibres de @@M5@@ sont géométriquement intègres, de même dimension @@M6@@. Soit @@M7@@ une famille de classes de faisceaux inversibles @@M8@@ sur les fibres de @@M9@@. Alors pour que @@M10@@ soit limitée, (il faut et) il suffit que dans les polynômes de Hilbert @@M11@@ des @@M12@@, tous les coefficients @@M13@@ restent bornés. De plus, si les fibres de @@M14@@ sont aussi géométriquement normales, alors il suffit que les trois coefficients @@M15@@ restent bornés.}

En effet, soit @@M16@@. Si @@M17@@, on peut remplacer @@M18@@ par la famille des classes des @@M19@@ grâce à 2.3 et 2.10\,; on peut ainsi supposer @@M20@@. Alors @@M21@@ est une @@M22@@-famille d'après 2.5, donc limitée si tous les coefficients @@M23@@ restent bornés, en vertu de 1.13, et, dans le cas où les fibres de @@M24@@ sont géométriquement normales, il suffit pour cela que les trois coefficients @@M25@@ restent bornés, en vertu de 2.9.

\subsection*{3. Théorèmes de finitude généraux}
\textbf{Lemme 3.1.} \emph{Soient @@M26@@ un sch\'ema et @@M27@@ un @@M28@@-sch\'ema quasi-compact et quasi-s\'epar\'e\,; soit @@M29@@ le foncteur de Picard, (localis\'e pour la topologie f.p.p.f.)}\SGAAuthorFootnote{*}{Cf. SGA 3 IV 6.3, et FGA n\textsuperscript{o} 232 pour la d\'efinition de @@M30@@, o\`u il convient ici, comme indiqu\'e, de remplacer la topologie fpqc par la topologie fppf.}.
\emph{Alors pour tout syst\`eme projectif filtrant @@M31@@ de @@M32@@-sch\'emas 
